Neste capítulo, são apresentados os experimentos realizados para avaliar o desempenho dos algoritmos Morpheus, Sunny e Rocket na geração automática de grades horárias acadêmicas do sistema IGRIS. O objetivo do benchmark é comparar o comportamento das três implementações diante do mesmo conjunto de testes, considerando métricas como tempo de execução, consumo de memória, número de gerações executadas, ocorrência de parada antecipada, score final e erro médio.

A comparação entre os algoritmos é necessária porque todos têm a mesma finalidade dentro do sistema: produzir uma grade horária acadêmica a partir das informações disponíveis sobre turmas, disciplinas, professores e horários. Entretanto, apesar de resolverem o mesmo problema, as implementações possuem características diferentes em relação à arquitetura, à organização interna dos dados, ao tratamento de conflitos, ao critério de parada e ao consumo de recursos computacionais.

Foram analisados dois cenários de execução, identificados nos relatórios como i3 e i5/full. Em ambos os casos, os algoritmos foram submetidos a execuções repetidas com variações de parâmetros evolutivos. Dessa forma, a análise não se limita a uma única execução isolada, mas considera o comportamento geral das implementações ao longo de diferentes configurações do algoritmo genético.

A comparação realizada neste capítulo não deve ser entendida como uma avaliação genérica entre linguagens de programação, arquiteturas ou tecnologias. O foco está nas três implementações concretas utilizadas no IGRIS. Portanto, os resultados discutidos dizem respeito ao comportamento observado no protocolo de benchmark executado, com os parâmetros, dados e ambiente disponíveis nos relatórios analisados.
